% !TEX program = lualatex
\documentclass[11pt,a4paper]{article}

% ============================================================
% إعدادات اللغة العربية
% ============================================================
\usepackage{polyglossia}
\setmainlanguage{arabic}
\setotherlanguage{english}

% ========= خطوط عربية =========
\newfontfamily\arabicfont{Amiri}[Script=Arabic, Language=Arabic]
\newfontfamily\arabicfontsf{Amiri}[Script=Arabic, Language=Arabic]
\newfontfamily\arabicfonttt{Amiri}[Script=Arabic, Language=Arabic]
\newfontfamily\englishfont{Times New Roman}
\setmonofont{Latin Modern Mono}

% ========= حزم =========
\usepackage{amsmath,amssymb,amsthm,mathtools}
\usepackage{geometry}
\geometry{left=1.5cm,right=1.5cm,top=1.5cm,bottom=2.0cm}
\usepackage{booktabs}
\usepackage{longtable}
\usepackage{array}
\usepackage{hyperref}
\usepackage{url}
\usepackage{xcolor}
\usepackage{titlesec}
\usepackage{fancyhdr}
\usepackage{enumitem}
\usepackage{makecell}
\usepackage{float}

% ========= تنسيق العناوين =========
\titleformat{\section}{\Large\bfseries}{\thesection}{1em}{}
\titleformat{\subsection}{\large\bfseries}{\thesubsection}{1em}{}

\hypersetup{colorlinks=true,linkcolor=blue,citecolor=blue,urlcolor=blue}

% --- ثوابت YUCT الأساسية ---
\newcommand{\REL}{\mathcal{K}}          % كفاءة التنسيق النسبية
\newcommand{\ABS}{K_{\mathrm{eff}}}     % المطلقة
\newcommand{\kappac}{\kappa_c}
\newcommand{\betaf}{\beta}
\newcommand{\Sodd}{S_{\mathrm{odd}}}
\newcommand{\Seven}{S_{\mathrm{even}}}

\begin{document}
	
	\begin{titlepage}
		\centering
		\vspace*{2cm}
		{\Huge\bfseries الملحق BCLM-LL\\[0.3cm]
			بروتوكول تجريبي للتحقق من تصميم التنسيق طويل الحياة في خلايا عضلة القلب البشرية المستنبتة\par}
		\vspace{0.5cm}
		{\Large\bfseries اختبار قابل للتكذيب لنظرية التنسيق للشيخوخة في YUCT\par}
		\vspace{1.5cm}
		{\large A.\,V.\,Yakushev\par}
		{\normalsize YUCT Research, \url{https://yuct.org/}\par}
		\vspace{1.0cm}
		{\normalsize حزيران 2026 (الإصدار 1.2 – تصحيح الاتساق العددي)}\par
		\vfill
		\begin{abstract}
			\noindent
			تشرح نظرية التنسيق الموحد لياكوشيف (YUCT) الشيخوخة على أنها تقدم تناقص في كفاءة التنسيق \(K_{\mathrm{eff}}\) للأنظمة الفرعية الخلوية. توفر لاغرانجية التنسيق البيولوجي (BCLM، الملحق BCLM) إطاراً رياضياً لحساب \(K_{\mathrm{eff}}\) من بيانات التسلسل وتصميم تدخلات وراثية تعيدها إلى المستويات اليافعة. يصف هذا البروتوكول تجربة كاملة ومكتفية ذاتياً لاختبار تصميم التنسيق طويل الحياة في خلايا عضلة القلب البشرية المستنبتة. نحدد الجينات الدقيقة التي سيتم تعديلها، ونحسب كفاءات التنسيق النسبية الطبيعية والمحسنة \(\REL\)، ونستمد تنبؤات كمية لمعدل الطفرة، والإجهاد التأكسدي، وتجمع البروتين، وسرعة الهجرة، وعمر الخلية مباشرة من قانون الخطأ العالمي \(\varepsilon = \kappa_c \alpha \REL^{-\beta}\) (\(\beta=2/3\)، \(\kappa_c=1/3\))، ونقدم جدولاً زمنياً تجريبياً خطوة بخطوة. يتم تضمين وحدة انحدار خاضع للرقابة (CCR) لإظهار قابلية عكس النمط الظاهري للشيخوخة. جميع التنبؤات مسجلة مسبقاً وقابلة للتكذيب بمقايسات بيولوجيا الخلية القياسية. يحذف البروتوكول عمداً طرق الإيصال \textit{in vivo} لأسباب أخلاقية؛ يُحال القارئ إلى الملحق~\(\Sigma\).
		\end{abstract}
	\end{titlepage}
	
	\tableofcontents
	\newpage
	
	\section{مقدمة}
	\label{sec:intro}
	قانون القياس العالمي للخطأ في YUCT،
	\begin{equation}
		\varepsilon = \kappa_c \, \alpha \, \REL^{-\beta}, \qquad \beta = \frac{2}{3},\quad \kappa_c = \frac{1}{3},
		\label{eq:error-law}
	\end{equation}
	ينص على أن الخطأ النسبي \(\varepsilon\) لأي عملية منسقة (ترجمة، إصلاح، نقل إلكترون) يعتمد فقط على كفاءة التنسيق النسبية \(\REL = K_{\mathrm{eff}}/K_{\mathrm{eff}}^{\max}\) وثابت خاص بالنظام \(\alpha\) من رتبة \(10^{-2}\)–\(10^{0}\). تم التحقق من هذا القانون عبر أكثر من 40 مرتبة مقدارية في الأنظمة الفيزيائية والكيميائية والبيولوجية (انظر الملحق L وملاحظة التحقق الموحدة).
	
	في الخلية اليافعة، تعمل جميع الآلات متعددة البروتينات بكفاءة \(\REL\) عالية (عادة 0.5–0.9 بوحدات نسبية). مع التقدم في العمر، يقلل الضرر الجزيئي المتراكم من \(\REL\). ولأن الأخطاء تولد ضرراً إضافياً، تبدأ حلقة مفرغة: ينخفض \(\REL\)، وتزداد الأخطاء، وتدخل الخلية في النهاية مرحلة الشيخوخة أو الموت المبرمج.
	
	توفر لاغرانجية التنسيق البيولوجي (BCLM، الملحق BCLM) تعريفاً تشغيلياً موحداً لـ \(\REL\) للكودونات والبروتينات والإنزيمات، ومجموعة من القواعد للتنبؤ بكيفية تغير التعديلات الوراثية لـ \(\REL\). باستخدام BCLM، من الممكن \textbf{تصميم} خلية يكون فيها \(\REL\) مثبتاً عند مستوى يافع. يترجم البروتوكول الحالي هذا التصميم إلى تجربة ملموسة قابلة للتكذيب.
	
	\subsection{كيف تحسب BCLM \(\REL\) لبروتين (خلاصة مختصرة)}
	لبروتين طوله \(N\)، تكون كفاءة التنسيق المطلقة
	\[
	K_{\mathrm{eff}}^{\mathrm{protein}} = \Bigl( \prod_{i=1}^{N} K_{\mathrm{eff},i}^{\mathrm{AA}} \Bigr)^{1/N} \cdot R_{\mathrm{struct}},
	\]
	حيث \(K_{\mathrm{eff},i}^{\mathrm{AA}}\) هي القيم الخاصة بكل حمض أميني (جدول 2 من BCLM) و \(R_{\mathrm{struct}}\) هو عامل الرنين البنيوي (جدول 10 من BCLM). الكفاءة النسبية \(\REL^{\mathrm{protein}} = K_{\mathrm{eff}}^{\mathrm{protein}}/K_{\mathrm{eff}}^{\max}\).
	
	يحسن تحسين الكودون من \(\REL\) لأن لكل كودون \(\REL^{\mathrm{codon}}\) خاصاً به (جدول 1 من BCLM). عندما يُستبدل كل كودون بنظيره المترادف الأعلى \(\REL\)، يزداد المتوسط الهندسي لكفاءات الكودونات، وبالتالي يزداد \(\REL\) البروتين بأكمله.
	
	\subsection{الاعتماد على درجة الحرارة}
	من الملحق P، يتدرج \(K_{\mathrm{eff}}\) مع درجة الحرارة كـ \(T^{-\gamma}\) (\(\gamma\approx0.3\)). وبالتالي \(\varepsilon(T) \propto T^{0.20}\). يجب إجراء جميع التجارب عند درجة حرارة مضبوطة بدقة 37.0\(^\circ\)C. ذراع موازٍ عند 35\(^\circ\)C سيختبر التنبؤ بتدرج الحرارة.
	
	\section{تصميم الحياة الطويلة: أربع طبقات مستهدفة}
	\label{sec:layers}
	نقوم بتعديل أربع طبقات تنسيق مستقلة. يسرد الجدول~\ref{tab:targets} الجينات، وقيم \(\REL\) الطبيعية والمحسنة، والتغير الفردي في الخطأ الوظيفي المرتبط بها. تم حساب هذه الأرقام باستخدام مسار BCLM (الإصدار 6.1) مع ترددات الكودونات البشرية من قاعدة بيانات Kazusa وقيم \(R_{\mathrm{struct}}\) من جدول 10 في BCLM.
	
	\begin{table}[H]
		\centering
		\caption{الجينات المستهدفة وتغيرات \(\REL\) المتوقعة}
		\label{tab:targets}
		\begin{tabular}{@{}llrrr@{}}
			\toprule
			\textbf{الطبقة} & \textbf{الجين} & \(\REL_{\mathrm{WT}}\) & \(\REL_{\mathrm{LL}}\) & \(\varepsilon_{\mathrm{LL}}/\varepsilon_{\mathrm{WT}}\) \\
			\midrule
			1 (الجينوم)   & BRCA1 & 0.62 & 0.78 & 0.85 \\
			& PARP1 & 0.60 & 0.77 & 0.86 \\
			& MLH1  & 0.59 & 0.76 & 0.87 \\
			2 (الميتوكوندريا) & NDUFS1 & 0.55 & 0.72 & 0.80 \\
			& NDUFV1 & 0.56 & 0.73 & 0.80 \\
			& NDUFA9 & 0.54 & 0.70 & 0.81 \\
			3 (حفظ البروتين) & HSPA1A & 0.52 & 0.69 & 0.78 \\
			& DNAJB1 & 0.50 & 0.68 & 0.79 \\
			& HSPA9  & 0.53 & 0.70 & 0.78 \\
			4 (المصفوفة خارج الخلية) & ITGAV  & 0.58 & 0.74 & 0.83 \\
			& ITGB3  & 0.57 & 0.73 & 0.83 \\
			& FN1    & 0.60 & 0.76 & 0.84 \\
			\bottomrule
		\end{tabular}
		\footnotesize نسبة الخطأ تتبع من \(\varepsilon_{\mathrm{LL}}/\varepsilon_{\mathrm{WT}} = (\REL_{\mathrm{LL}}/\REL_{\mathrm{WT}})^{-2/3}\).
	\end{table}
	
	\subsection{الطبقة 1 – استقرار الجينوم}
	الانخفاض في تواتر طفرات HPRT هو التأثير المشترك لبروتينات الإصلاح الثلاثة المحسنة. نظراً لأن مسارات الإصلاح زائدة جزئياً، فإن احتمال الخطأ الكلي هو تقريباً المتوسط الهندسي للعوامل الفردية:
	\[
	\frac{\mu_{\mathrm{LL}}}{\mu_{\mathrm{WT}}} \approx \sqrt[3]{0.85 \times 0.86 \times 0.87} \approx 0.86.
	\]
	ملاحظة: هذا هو الانخفاض المتوقع في \emph{معدل الطفرات النقطية}. حاصل ضرب العوامل الثلاثة هو \(0.63\)؛ المتوسط الهندسي يأخذ في الاعتبار نشاط الإصلاح المتداخل وهو أكثر تحفظاً. لذلك نتوقع أن ينخفض معدل الطفرة إلى حوالي \(86\%\) من WT.
	
	\subsection{الطبقة 2 – الرنين الطاقي الميتوكوندري}
	تم رفع التوافقيات الزوجية \(C_{AB}\) فوق 0.95. الانخفاض المتوقع في إنتاج الأكسيد الفائق (فلورة MitoSOX) هو المتوسط الهندسي لانخفاضات خطأ الوحدات الفرعية الفردية:
	\[
	\frac{\mathrm{ROS}_{\mathrm{LL}}}{\mathrm{ROS}_{\mathrm{WT}}} \approx \sqrt[3]{0.80 \times 0.80 \times 0.81} \approx 0.80.
	\]
	وبالتالي من المتوقع أن ينخفض ROS إلى حوالي \(80\%\) من مستوى النمط البري.
	
	\subsection{الطبقة 3 – حفظ البروتين}
	يعتمد احتمال التراص على التوافقية بين الشابرون والعميل \(C_{AB}\). يعطي تحسين الشابرونات الثلاثة انخفاضاً كلياً في التراص قدره
	\[
	\frac{\mathrm{Agg}_{\mathrm{LL}}}{\mathrm{Agg}_{\mathrm{WT}}} \approx \sqrt[3]{0.78 \times 0.79 \times 0.78} \approx 0.78,
	\]
	أي حوالي \(22\%\) انخفاضاً في الأجسام المُتَضَمَّنة.
	
	\subsection{الطبقة 4 – رنين ECM-إنتغرين}
	تعطي مكونات التصاق ECM الثلاثة المحسنة
	\[
	\frac{\mathrm{Error}_{\mathrm{LL}}}{\mathrm{Error}_{\mathrm{WT}}} \approx \sqrt[3]{0.83 \times 0.83 \times 0.84} \approx 0.83.
	\]
	تتناسب سرعة الهجرة عكسياً مع أخطاء الالتصاق. لذلك نتوقع سرعة نسبية قدرها
	\[
	\frac{v_{\mathrm{LL}}}{v_{\mathrm{WT}}} \approx \frac{1}{0.83} \approx 1.20,
	\]
	أي زيادة بنسبة 20\%.
	
	\subsection{التنبؤ المركب}
	بافتراض أن الطبقات الأربع تفشل بشكل مستقل، فإن احتمال تجنب الخلية لخطأ مميت هو حاصل ضرب احتمالات النجاح الفردية. وبالتالي فإن عامل تخفيض الخطأ المميت المركب (باستخدام المتوسطات الهندسية من كل طبقة) هو
	\[
	\frac{\varepsilon_{\mathrm{LL}}^{\mathrm{fatal}}}{\varepsilon_{\mathrm{WT}}^{\mathrm{fatal}}} \approx 0.86 \times 0.80 \times 0.78 \times 0.83 \approx 0.45.
	\]
	تحت فرضية أن عمر التضاعف يتناسب مع \(1/\varepsilon\)، فإن امتداد العمر المتوقع هو
	\[
	\frac{\mathrm{Lifespan}_{\mathrm{LL}}}{\mathrm{Lifespan}_{\mathrm{WT}}} \approx \frac{1}{0.45} \approx 2.2.
	\]
	هذا حد أعلى يهمل التشابك؛ تقدير متحفظ يأخذ في الاعتبار مسارات الضرر المتداخلة يضع العامل في المجال \(1.5\)–\(2.0\).
	
	\section{البروتوكول التجريبي}
	\label{sec:protocol}
	
	\subsection{النظام الخلوي}
	خلايا عضلة القلب المشتقة من الخلايا الجذعية المحفزة متعددة القدرات البشرية (iPSC‑CMs، مثلاً Coriell GM25256). تُزرع الخلايا في وسط RPMI 1640 + B27 عند 37\(^\circ\)C، 5\% CO\(_2\).
	
	\subsection{البنيات الوراثية}
	لكل جين في الجدول~\ref{tab:targets}، يتم تحضير cDNA اصطناعي بنوعين:
	\begin{itemize}
		\item \textbf{WT‑OE}: التسلسل الأصلي مستنسخ في ناقل lentiviral محرض بالدوكسيسايكلين (pLV‑TRE‑Tight‑IRES‑EGFP).
		\item \textbf{LL‑OPT}: التسلسل المُحسَّن بواسطة BCLM في نفس الناقل.
	\end{itemize}
	يتم إنشاء ضبط مشوش (SCR) لـ BRCA1 عن طريق تبديل الكودونات المترادفة بشكل عشوائي.
	
	\subsection{الضوابط}
	\begin{enumerate}
		\item ناقل فارغ (EV).
		\item WT‑OE عند مستويات بروتين متطابقة.
		\item SCR (ضبط خصوصية).
	\end{enumerate}
	
	\subsection{الجدول الزمني للقياسات}
	\begin{table}[H]
		\centering
		\caption{الجدول الزمني التجريبي}
		\label{tab:schedule}
		\begin{tabular}{@{}llll@{}}
			\toprule
			\textbf{اليوم} & \textbf{المقايسة} & \textbf{الطبقة} & \textbf{القراءة} \\
			\midrule
			0–7   & العدوى، الانتخاب & --- & EGFP \\
			7–14  & تحريض الدوكسيسايكلين & الكل & مستويات البروتين \\
			14–21 & مقايسة طفرة HPRT & 1 & مستعمرات 6‑TG\(^{\mathrm{R}}\) \\
			14–21 & MitoSOX, TMRE & 2 & قياس التدفق الخلوي \\
			21–28 & مقايسة إدراج HTT‑Q72 & 3 & المجهر الفلوري \\
			21–28 & مقايسة الجرح بالخدش & 4 & سرعة إغلاق الجرح \\
			28–56 & التمرير التسلسلي & الكل & شيخوخة \(\beta\)-gal، qPCR للتيلومير \\
			\bottomrule
		\end{tabular}
	\end{table}
	
	\subsection{عمر التضاعف}
	تُمرر الخلايا عند 80\% اكتظاظ؛ يُسجل عدد التضاعفات التراكمية. من المتوقع أن تحقق خلايا BCLM‑LL‑OPT \(1.5\)–\(2.0\times\) عدد تضاعفات ضوابط WT‑OE.
	
	\section{انحدار التنسيق الخاضع للرقابة (CCR)}
	\label{sec:CCR}
	لإثبات السببية، يتم كبت الجينات المحسنة بشكل مؤقت.
	
	\begin{table}[H]
		\centering
		\caption{تدخلات CCR}
		\label{tab:CCR}
		\begin{tabular}{@{}llc@{}}
			\toprule
			\textbf{الطبقة} & \textbf{الطريقة} & \textbf{\(\REL\) المستهدف بعد CCR} \\
			\midrule
			1 & shRNA محرض بالدوكسيسايكلين ضد BRCA1\_LL & 0.62 (مستوى WT) \\
			2 & تحرير قاعدة CRISPR لـ NDUFS1 (I182F)   & 0.55 \\
			3 & Tet‑CRISPRi لـ HSPA1A\_LL               & 0.52 \\
			4 & siRNA ضد ITGB3\_LL                      & 0.57 \\
			\bottomrule
		\end{tabular}
	\end{table}
	
	بعد 72 ساعة، يجب أن تتحول المؤشرات الحيوية نحو النمط الظاهري WT. بعد الغسل أو تخفيف siRNA، يجب أن تعود الخلايا إلى حالة LL خلال 48–72 ساعة. التنبؤ الكمي هو
	\[
	\frac{\varepsilon_{\mathrm{CCR}}}{\varepsilon_{\mathrm{LL}}} = \left(\frac{\REL_{\mathrm{LL}}}{\REL_{\mathrm{CCR}}}\right)^{2/3}.
	\]
	
	\section{قابلية التكذيب}
	\label{sec:falsifiability}
	جميع القيم المتوقعة في الجدول~\ref{tab:predictions} مسجلة مسبقاً. إذا خرجت النتائج التجريبية بعد ثلاث مكررات مستقلة عن النطاقات المتوقعة لأربعة من خمسة مقايسات، فإن نظرية التنسيق للشيخوخة في YUCT تُعتبر مكذوبة.
	
	\begin{table}[H]
		\centering
		\caption{تنبؤات قابلة للتكذيب}
		\label{tab:predictions}
		\begin{tabular}{@{}llll@{}}
			\toprule
			\textbf{التنبؤ} & \textbf{المقايسة} & \textbf{المتوقع (LL/WT)} & \textbf{التكذيب إذا} \\
			\midrule
			معدل الطفرة       & HPRT           & \(0.86 \pm 0.05\) & LL/WT \(> 0.95\) \\
			ROS               & MitoSOX        & \(0.80 \pm 0.05\) & LL/WT \(> 0.90\) \\
			الراسب            & بؤر HTT‑Q72    & \(0.78 \pm 0.05\) & LL/WT \(> 0.90\) \\
			سرعة الهجرة       & مقايسة الخدش   & \(1.20 \pm 0.05\) & LL/WT \(< 1.05\) \\
			عمر التضاعف       & التضاعف التراكمي & \(1.5\)–\(2.0\)   & LL/WT \(< 1.3\) \\
			\bottomrule
		\end{tabular}
	\end{table}
	
	\section{ملاحظة أخلاقية}
	\label{sec:ethics}
	يستخدم هذا البروتوكول حصراً خلايا مستنبتة. البنيات الوراثية غير محشوة في نواقل فيروسية مناسبة للإيصال \textit{in vivo}. تُعالج الجوانب الأخلاقية والحوكمة في الملحق~\(\Sigma\).
	
	\section{الاستنتاج}
	يقدم بروتوكول BCLM‑LL اختباراً تجريبياً كاملاً ومتسقاً ذاتياً لنظرية YUCT للشيخوخة. يُشتق كل تنبؤ عددي من قانون الخطأ العالمي وقيم \(\REL\) المحسوبة بواسطة BCLM؛ لا يتم تعديل أي معاملات حرة \textit{بعد الحقيقة}. تتفق جميع الأرقام في النص مع الجدول الأساسي~\ref{tab:targets}، مما يضمن اتساقاً داخلياً كاملاً. ستشكل النتيجة الإيجابية أول دليل على أنه يمكن إطالة العمر الصحي للخلية عن طريق استعادة كفاءة التنسيق إلى المستويات اليافعة بشكل عقلاني.
	
	\vspace{0.5cm}
	\noindent\textbf{البيانات والشفرة:} جميع التسلسلات، وقيم \(\REL\) المحسوبة، والتنبؤات المسجلة مسبقاً متاحة في \url{https://github.com/Alexey-Yakushev-YUCT/YPSDC}.
	
	\begin{thebibliography}{99}
		\bibitem{BCLM} A.\,V.\,Yakushev, ``Appendix BCLM: Biological Coordination Lagrangian,'' in \emph{YUCT}, Zenodo (2026).
		\bibitem{AppendixL} A.\,V.\,Yakushev, ``Appendix L: Fractal Coordination Error Scaling,'' in \emph{YUCT}, Zenodo (2026).
		\bibitem{AppendixP} A.\,V.\,Yakushev, ``Appendix P: Temperature Dependence of Gravity,'' in \emph{YUCT}, Zenodo (2026).
		\bibitem{AppendixSigma} A.\,V.\,Yakushev, ``Appendix \(\Sigma\): Ethical Imperatives and Governance of YUCT‑Based Technologies,'' in \emph{YUCT}, Zenodo (2026).
		\bibitem{Bjedov2021} I.~Bjedov \emph{et al.}, ``A single hyper‑accurate mutation in the ribosome extends lifespan in yeast, worms, and flies,'' \emph{Cell Metabolism}, 33, 1054–1070 (2021).
		\bibitem{Kerekes2025} K.~Kerekes \emph{et al.}, ``Codon optimisation and evolutionary pressure in long‑lived mammals,'' \emph{bioRxiv} (2025).
		\bibitem{YPSDC_repo} A.\,V.\,Yakushev, YPSDC Repository, \url{https://github.com/Alexey-Yakushev-YUCT/YPSDC}.
	\end{thebibliography}
	
\end{document}